
\begin{unnumbered}

\subsection{1 НАЙМЕНУВАННЯ ТА ОБЛАСТЬ ЗАСТОСУВАННЯ}

Це технічне завдання стосується розробки кросплатформного мобільного застосунку для відстеження розвитку немовляти. Застосунок призначений для батьків та опікунів, які потребують зручного інструменту для моніторингу активностей дитини, контролю антропометричних показників відповідно до стандартів ВООЗ, управління вакцинаціями та досягненнями, а також отримання персоналізованих педіатричних рекомендацій. Система розрахована на використання в побутових умовах на мобільних пристроях під управлінням iOS та Android.

\subsection{2 ПІДСТАВИ ДЛЯ РОЗРОБКИ}

Підставою для розробки цього програмного продукту є виконання вимог кваліфікаційно-освітнього рівня «бакалавр» зі спеціальності «Комп'ютерна інженерія», затвердженого відповідними освітніми та науковими установами Національного технічного університету України «Київський політехнічний інститут імені Ігоря Сікорського».

\subsection{3 МЕТА ТА ПРИЗНАЧЕННЯ РОЗРОБКИ}

Метою розробки є створення зручного та функціонального кросплатформного мобільного застосунку, який допоможе батькам систематично відстежувати розвиток немовляти в перші роки життя.

Застосунок призначений для ведення щоденника активностей дитини (годування, сон, підгузки), моніторингу росту та ваги з порівнянням до стандартів ВООЗ, а також відстеження вакцинацій і вікових досягнень. Система забезпечує синхронізацію даних між пристроями в реальному часі, що дозволяє обом батькам мати актуальну інформацію одночасно.

\subsection{4 ДЖЕРЕЛА РОЗРОБКИ}

Джерелами розробки даного дипломного проєкту є офіційна технічна документація до використовуваних технологій (React Native, Expo, Supabase), стандарти ВООЗ щодо показників росту та розвитку дітей, наукові публікації у сфері мобільної розробки, а також статті та матеріали в мережі Інтернет за тематикою проєкту.

\subsection{5 ТЕХНІЧНІ ВИМОГИ}

\subsubsection{5.1. Вимоги до розробленого продукту}

Розроблена система повинна задовольняти наступні вимоги:

\begin{itemize}
  \item кросплатформна підтримка iOS 16.0 і новіших та Android 10 (API~29) і новіших версій;
  \item автентифікація користувачів та ізоляція даних на рівні рядків бази даних (Row-Level Security);
  \item синхронізація даних між пристроями в реальному часі через Supabase Realtime;
  \item відображення графіків зростання та порівняння з еталонними кривими ВООЗ;
  \item підтримка кількох профілів дітей в рамках одного облікового запису;
  \item зрозумілий та інтуїтивний інтерфейс, що не вимагає попередньої технічної підготовки.
\end{itemize}

\subsubsection{5.2. Вимоги до програмного забезпечення}

\begin{itemize}
  \item операційна система: iOS 16.0+ або Android 10+;
  \item активне підключення до мережі Інтернет для синхронізації даних між пристроями;
  \item для розробки: Node.js 18+, Expo CLI, Git.
\end{itemize}

\subsubsection{5.3. Вимоги до апаратної частини}

\begin{itemize}
  \item смартфон з оперативною пам'яттю не менше 3~ГБ;
  \item наявність підключення до мережі Інтернет (Wi-Fi або мобільний інтернет).
\end{itemize}

\subsection{6 ЕТАПИ РОЗРОБКИ}

\begin{table}[htbp]
\begin{center}
\begin{tabular}{|p{10cm}|p{5cm}|}
\hline
\textbf{Назва етапу виконання} & \textbf{Термін виконання} \\
\hline
Затвердження теми роботи & 06.04.2026--07.06.2026 \\
\hline
Вивчення та аналіз завдання & 02.05.2026--08.05.2026 \\
\hline
Розробка архітектури та загальної структури системи & 09.05.2026--15.05.2026 \\
\hline
Розробка структур окремих частин системи & 09.05.2026--15.05.2026 \\
\hline
Програмна реалізація системи & 16.05.2026--22.05.2026 \\
\hline
Виправлення помилок & 23.05.2026--29.05.2026 \\
\hline
Оформлення пояснювальної записки & 30.05.2026--07.06.2026 \\
\hline
\end{tabular}
\end{center}
\end{table}

\end{unnumbered}
